    % ==== Deep Learning =============================================================================
    \chapter{Deep Learning}

% ==== Section 1: Introduction to Neural Networks and Deep Learning ===========================
\section{Introduction to Neural Networks and Deep Learning}

% ---- 6.1.1 The Biological Brain -------------------------------------------------------------
\subsection{The Biological Brain}

\begin{tcbitemize}[ skin=sectionraster ]
    \tcbitem[title={The Biological Neuron}, raster multicolumn=3]
    Artificial neural networks draw inspiration from the structure of biological neurons\textsuperscript{\cite[][pp.~27--35]{Kriesel2007NeuralNetworks}}.
    A biological neuron consists of:
    \begin{itemize}
        \item \textbf{Dendrites}: tree-like input branches that receive signals from other neurons
        \item \textbf{Soma} (cell body): integrates incoming signals
        \item \textbf{Axon}: long output fibre that transmits the signal onward
        \item \textbf{Synapses}: junctions between neurons; neurotransmitter release modulates connection strength
    \end{itemize}

    \tcbitem[title={From Biology to Computation}, raster multicolumn=3]
    The neuron's membrane rests at approximately $-70\,\text{mV}$.
    When excitatory inputs push the potential above a threshold ($\approx -55\,\text{mV}$), an \textbf{action potential} (spike) propagates down the axon\textsuperscript{\cite[][pp.~29--31]{Kriesel2007NeuralNetworks}}.
    Chemical synapses strengthen or weaken over time (\textbf{synaptic plasticity}), providing a biological basis for learning.
    \tcblower
    The \textbf{100-step rule}: the brain performs complex recognition in ${\sim}500\,\text{ms}$; neurons fire at ${\sim}100\,\text{Hz}$, allowing at most ${\sim}100$ sequential processing steps.
    This implies massive parallelism --- motivating parallel architectures in artificial networks.
\end{tcbitemize}

% ---- 6.1.2 Perceptron and Multi-Layer Perceptrons -------------------------------------------
\subsection{Perceptron and Multi-Layer Perceptrons}

\begin{tcbitemize}[ skin=sectionraster ]
    \tcbitem[title={The Perceptron}, raster multicolumn=3]
    The \textbf{perceptron} (Rosenblatt, 1957) is the simplest artificial neuron\textsuperscript{\cite[][pp.~57--62]{Kriesel2007NeuralNetworks}}.
    It computes a weighted sum of inputs and applies a binary threshold:
    \begin{equation}
        y = \begin{cases} 1 & \text{if } \mathbf{w}^\top \mathbf{x} + b \geq 0 \\ 0 & \text{otherwise} \end{cases}
    \end{equation}
    The \textbf{perceptron convergence theorem} guarantees that the learning rule will find a separating hyperplane if the data is linearly separable.

    \tcbitem[title={The XOR Problem and MLPs}, raster multicolumn=3]
    Minsky \& Papert (1969) proved that a single perceptron \emph{cannot} solve the XOR problem --- it can only represent linearly separable functions\textsuperscript{\cite[][pp.~63--65]{Kriesel2007NeuralNetworks}}.
    This limitation contributed to the first ``AI winter''.
    The solution is the \textbf{Multi-Layer Perceptron (MLP)}: stacking layers with nonlinear activation functions creates networks that can represent arbitrarily complex decision boundaries.
    \tcblower
    \begin{center}
    \begin{tikzpicture}[
        node distance=0.8cm and 1.2cm,
        neuron/.style={circle, draw, minimum size=0.6cm, font=\scriptsize, inner sep=1pt},
        arr/.style={-{Stealth[length=1.5mm]}, thick}
    ]
        % Input layer
        \node[neuron, fill=blue!15] (i1) {$x_1$};
        \node[neuron, fill=blue!15, below=of i1] (i2) {$x_2$};
        % Hidden layer
        \node[neuron, fill=orange!20, right=of i1, yshift=-0.4cm] (h1) {$h_1$};
        \node[neuron, fill=orange!20, below=of h1] (h2) {$h_2$};
        % Output
        \node[neuron, fill=green!20, right=of h1, yshift=-0.4cm] (o) {$y$};

        \foreach \i in {i1,i2} \foreach \h in {h1,h2} \draw[arr] (\i) -- (\h);
        \foreach \h in {h1,h2} \draw[arr] (\h) -- (o);

        \node[below=0.15cm of i2, font=\sffamily\tiny] {Input};
        \node[below=0.15cm of h2, font=\sffamily\tiny] {Hidden};
        \node[below=0.55cm of o, font=\sffamily\tiny] {Output};
    \end{tikzpicture}
\captionof{figure}{Perceptron and Multi-Layer Perceptrons: The XOR Problem and MLPs.}
\label{fig:ch06-deep-learning-inline-01}
    \end{center}

    \tcbitem[title={From Perceptrons to Deep Learning}, raster multicolumn=6]
    The transition from single perceptrons to deep networks involved three key advances\textsuperscript{\cite[][pp.~1--3]{fleuretLittleBookDeep}}:
    \begin{itemize}
        \item \textbf{Hebbian learning} $\to$ \textbf{error-correction rules} $\to$ \textbf{backpropagation} (Rumelhart, Hinton \& Williams, 1986): enabled training of multi-layer networks by propagating gradients backward through layers
        \item \textbf{Nonlinear activations}: sigmoid, tanh, and later ReLU allow hidden layers to learn complex feature hierarchies
        \item \textbf{Hardware and data}: GPUs and large datasets made training deep networks practical from $\sim$2012 onward
    \end{itemize}
    \textbf{Deep learning} refers to neural networks with many hidden layers that learn hierarchical representations --- from low-level features (edges, textures) to high-level concepts (objects, semantics).
\end{tcbitemize}

% ---- 6.1.3 Activation Functions -------------------------------------------------------------
\subsection{Activation Functions}

\begin{tcbitemize}[ skin=sectionraster ]
    \tcbitem[title={Why Activation Functions Matter}, raster multicolumn=6]
    Without nonlinear activation functions, a multi-layer network collapses to a single linear transformation\textsuperscript{\cite[][p.~17]{fleuretLittleBookDeep}}.
    The choice of activation function affects gradient flow, training speed, and representational capacity.
    \tcblower
    \begin{tblr}{|Q[l,m]|X[l]|X[l]|X[l]|}
        \hline
        \textbf{Function} & \textbf{Formula} & \textbf{Range} & \textbf{Key Property} \\
        \hline
        Sigmoid & $\sigma(x) = \frac{1}{1+e^{-x}}$ & $(0, 1)$ & Vanishing gradients for $|x| \gg 0$ \\
        \hline
        Tanh & $\tanh(x)$ & $(-1, 1)$ & Zero-centred; still saturates \\
        \hline
        ReLU & $\max(0, x)$ & $[0, \infty)$ & Fast; can ``die'' (zero gradient for $x < 0$) \\
        \hline
        Leaky ReLU & $\max(\alpha x, x)$ & $(-\infty, \infty)$ & Prevents dying neurons ($\alpha \approx 0.01$) \\
        \hline
        GELU & $x \cdot \Phi(x)$ & $\approx (-0.17, \infty)$ & Smooth; used in Transformers \\
        \hline
        Softmax & $\frac{e^{x_i}}{\sum_j e^{x_j}}$ & $(0, 1)$, sums to 1 & Output layer for classification \\
        \hline
    \end{tblr}
\captionof{table}{Activation Functions: Why Activation Functions Matter.}
\label{tab:ch06-deep-learning-inline-01}

    \tcbitem[title={ReLU: The Default Choice}, raster multicolumn=3]
    The \textbf{Rectified Linear Unit} $\mathrm{ReLU}(x) = \max(0, x)$ became the default activation after Glorot et al.\ (2011) showed it greatly improves training of deep networks\textsuperscript{\cite[][]{glorotUnderstandingDifficultyTraining}}.
    Its gradient is either $0$ or $1$ --- avoiding the vanishing gradient problem that plagues sigmoid and tanh in deep architectures.
    \tcblower
    The \textbf{dying ReLU problem}: if a neuron's pre-activation is always negative, its gradient is permanently zero and it stops learning.
    Leaky ReLU and parametric ReLU (PReLU) mitigate this by allowing a small gradient for negative inputs.

    \tcbitem[title={Softmax for Classification}, raster multicolumn=3]
    For multi-class classification, the \textbf{softmax} function converts a vector of raw scores (logits) into a probability distribution\textsuperscript{\cite[][p.~18]{fleuretLittleBookDeep}}:
    \begin{equation}
        \mathrm{softmax}(x_i) = \frac{e^{x_i}}{\sum_{j=1}^{K} e^{x_j}}
    \end{equation}
    All outputs are positive and sum to $1$.
    Combined with cross-entropy loss, softmax provides useful gradients for training classifiers. Its outputs sum to one, but they are not automatically calibrated probabilities; calibration must be evaluated on held-out data.
\end{tcbitemize}


% ==== Section 2: Network Architectures =======================================================
\section{Network Architectures}

% ---- 6.2.1 Feed-Forward Networks ------------------------------------------------------------
\subsection{Feed-Forward Networks}

\begin{tcbitemize}[ skin=sectionraster ]
    \tcbitem[title={Dense Layers}, raster multicolumn=3]
    A \textbf{feed-forward network} (also called a \emph{fully connected} or \emph{dense} network) consists of layers where every neuron connects to all neurons in the next layer\textsuperscript{\cite[][pp.~15--17]{fleuretLittleBookDeep}}.
    Each layer computes:
    \begin{equation}
        Y = \sigma(WX + b)
    \end{equation}
    where $W$ is the weight matrix, $b$ the bias vector, and $\sigma$ a nonlinear activation.
    Information flows strictly forward --- there are no cycles or feedback connections.

    \tcbitem[title={Universal Approximation}, raster multicolumn=3]
    The \textbf{Universal Approximation Theorem} (Cybenko, 1989; Hornik, 1991) states that a feed-forward network with a single hidden layer containing sufficiently many neurons can approximate any continuous function on a compact set to arbitrary precision\textsuperscript{\cite[][p.~16]{fleuretLittleBookDeep}}.
    However, ``sufficiently many'' may be exponentially large.
    \textbf{Deeper networks} are exponentially more parameter-efficient than wider shallow ones for many function classes\textsuperscript{\cite[][pp.~8--12]{robertsPrinciplesDeepLearning2022}} --- this is the fundamental motivation for \emph{deep} learning.
\end{tcbitemize}

% ---- 6.2.2 Convolutional Networks -----------------------------------------------------------
\subsection{Convolutional Networks}

\begin{tcbitemize}[ skin=sectionraster ]
    \tcbitem[title={The Convolution Operation}, raster multicolumn=3]
    A \textbf{Convolutional Neural Network (CNN)} replaces dense matrix multiplications with convolution operations\textsuperscript{\cite[][pp.~43--48]{fleuretLittleBookDeep}}.
    Small learned \textbf{filters} (kernels) slide across the input, computing dot products at each position.
    Key properties:
    \begin{itemize}
        \item \textbf{Weight sharing}: the same filter is applied everywhere $\to$ translation equivariance
        \item \textbf{Local connectivity}: each output depends only on a small input region
        \item \textbf{Parameter efficiency}: far fewer parameters than a dense layer of the same input/output size
    \end{itemize}

    \tcbitem[title={Pooling and Feature Hierarchies}, raster multicolumn=3]
    \textbf{Pooling layers} (max or average) reduce spatial dimensions and provide a degree of translation invariance\textsuperscript{\cite[][pp.~49--50]{fleuretLittleBookDeep}}.
    Stacking convolution and pooling layers creates a \textbf{feature hierarchy}: early layers detect edges and textures; deeper layers detect parts and objects.
    The \textbf{receptive field} --- the input region that influences a single output unit --- grows with network depth.

    \tcbitem[title={Landmark CNN Architectures}, raster multicolumn=6]
    \tcblower
    \begin{tblr}{|Q[l,m]|Q[c,m]|X[l]|}
        \hline
        \textbf{Architecture} & \textbf{Year} & \textbf{Key Innovation} \\
        \hline
        LeNet-5 & 1998 & First successful CNN; handwritten digit recognition \\
        \hline
        AlexNet & 2012 & ReLU activation, dropout, GPU training; won ImageNet \\
        \hline
        VGG & 2014 & Uniform $3\times3$ filters; demonstrated depth matters \\
        \hline
        ResNet & 2015 & Skip connections; enabled training of 100+ layer networks \\
        \hline
    \end{tblr}
\captionof{table}{Convolutional Networks: Landmark CNN Architectures.}
\label{tab:ch06-deep-learning-inline-02}
\end{tcbitemize}

% ---- 6.2.3 Recurrent Networks, Memory Cells and LSTMs --------------------------------------
\subsection{Recurrent Networks, Memory Cells and LSTMs}

\begin{tcbitemize}[ skin=sectionraster ]
    \tcbitem[title={Recurrent Neural Networks}, raster multicolumn=3]
    A \textbf{Recurrent Neural Network (RNN)} processes sequential data by maintaining a hidden state $h_t$ that is updated at each time step\textsuperscript{\cite[][pp.~55--58]{fleuretLittleBookDeep}}:
    \begin{equation}
        h_t = \sigma(W_h h_{t-1} + W_x x_t + b)
    \end{equation}
    Weight sharing across time steps means the same parameters handle sequences of any length.
    The \textbf{vanishing gradient problem} makes it difficult for standard RNNs to learn long-range dependencies: gradients are multiplied through many time steps and either vanish or explode.

    \tcbitem[title={LSTM: Long Short-Term Memory}, raster multicolumn=3]
    The \textbf{LSTM} (Hochreiter \& Schmidhuber, 1997) solves the vanishing gradient problem with a \textbf{cell state} $c_t$ and three multiplicative gates\textsuperscript{\cite[][pp.~59--61]{fleuretLittleBookDeep}}:
    \begin{itemize}
        \item \textbf{Forget gate} $f_t$: decides what to discard from cell state
        \item \textbf{Input gate} $i_t$: decides what new information to store
        \item \textbf{Output gate} $o_t$: decides what to output from cell state
    \end{itemize}
    \tcblower
    \begin{equation}
        c_t = f_t \odot c_{t-1} + i_t \odot \tilde{c}_t, \quad h_t = o_t \odot \tanh(c_t)
    \end{equation}
    The cell state provides a ``highway'' for gradients to flow through many time steps unimpeded.

    \tcbitem[title={GRU: Gated Recurrent Unit}, raster multicolumn=6]
    The \textbf{GRU} (Cho et al., 2014) simplifies the LSTM by merging the cell state and hidden state and using only two gates\textsuperscript{\cite[][p.~62]{fleuretLittleBookDeep}}:
    a \emph{reset gate} $r_t$ (controls how much past state to forget) and an \emph{update gate} $z_t$ (controls the mix of old state and new candidate).
    GRUs have fewer parameters than LSTMs and often achieve similar performance on many tasks.
    Both architectures have been largely superseded by Transformer-based models for most sequence tasks, but remain relevant for streaming and low-latency applications.
\end{tcbitemize}


% ==== Section 3: Neural Network Training =====================================================
\section{Neural Network Training}

% ---- 6.3.1 Forward Pass and Loss Functions --------------------------------------------------
\subsection{Forward Pass and Loss Functions}

\begin{tcbitemize}[ skin=sectionraster ]
    \tcbitem[title={The Forward Pass}, raster multicolumn=3]
    During the \textbf{forward pass}, input data propagates through the network layer by layer\textsuperscript{\cite[][pp.~20--22]{fleuretLittleBookDeep}}.
    Each layer applies its learned transformation (weights, bias, activation), producing intermediate representations.
    The final layer produces a prediction $\hat{y}$ that is compared to the true target $y$ using a \textbf{loss function} $\mathcal{L}(y, \hat{y})$.

    \tcbitem[title={Common Loss Functions}, raster multicolumn=3]
    \tcblower
    \begin{tblr}{|Q[l,m]|X[l]|Q[l,m]|}
        \hline
        \textbf{Loss} & \textbf{Formula} & \textbf{Task} \\
        \hline
        MSE & $\frac{1}{N}\sum(y_i - \hat{y}_i)^2$ & Regression \\
        \hline
        MAE & $\frac{1}{N}\sum|y_i - \hat{y}_i|$ & Regression (robust) \\
        \hline
        Cross-entropy & $-\sum y_i \log \hat{y}_i$ & Classification \\
        \hline
        Binary CE & $-[y\log\hat{y} + (1{-}y)\log(1{-}\hat{y})]$ & Binary classif. \\
        \hline
    \end{tblr}
\captionof{table}{Forward Pass and Loss Functions: Common Loss Functions.}
\label{tab:ch06-deep-learning-inline-03}
\end{tcbitemize}

% ---- 6.3.2 Weight Initialization and Transfer Function --------------------------------------
\subsection{Weight Initialization and Transfer Function}

\begin{tcbitemize}[ skin=sectionraster ]
    \tcbitem[title={Why Initialization Matters}, raster multicolumn=3]
    Initializing all weights to zero causes a \textbf{symmetry problem}: all neurons in a layer compute the same output and receive the same gradient, so they never differentiate\textsuperscript{\cite[][pp.~8--10]{robertsPrinciplesDeepLearning2022}}.
    Too-large initial weights cause activations and gradients to explode; too-small weights cause them to vanish.
    Proper initialization keeps the variance of activations and gradients approximately constant across layers.

    \tcbitem[title={Xavier and He Initialization}, raster multicolumn=3]
    \textbf{Xavier/Glorot initialization}\textsuperscript{\cite[][]{glorotUnderstandingDifficultyTraining}} (for sigmoid/tanh):
    \begin{equation}
        W \sim \mathcal{N}\!\left(0,\; \frac{2}{n_{\mathrm{in}} + n_{\mathrm{out}}}\right)
    \end{equation}
    \textbf{He initialization} (for ReLU):
    \begin{equation}
        W \sim \mathcal{N}\!\left(0,\; \frac{2}{n_{\mathrm{in}}}\right)
    \end{equation}
    These schemes ensure that the expected variance of the output of each layer matches the variance of its input, preventing signal degradation in deep networks.
\end{tcbitemize}

% ---- 6.3.3 Backpropagation and Gradient Descent ---------------------------------------------
\subsection{Backpropagation and Gradient Descent}

\begin{tcbitemize}[ skin=sectionraster ]
    \tcbitem[title={Backpropagation}, raster multicolumn=3]
    \textbf{Backpropagation} computes the gradient of the loss with respect to every parameter by applying the chain rule layer by layer, from output back to input\textsuperscript{\cite[][pp.~1--4]{lecunTheoreticalFrameworkBack}}.
    For a network with layers $f_1, f_2, \ldots, f_L$:
    \begin{equation}
        \frac{\partial \mathcal{L}}{\partial W_l} = \frac{\partial \mathcal{L}}{\partial a_L} \cdot \frac{\partial a_L}{\partial a_{L-1}} \cdots \frac{\partial a_{l+1}}{\partial a_l} \cdot \frac{\partial a_l}{\partial W_l}
    \end{equation}
    Each layer's local Jacobian is computed once and reused --- making the algorithm $O(N)$ in the number of parameters.

    \tcbitem[title={Gradient Descent Variants}, raster multicolumn=3]
    Backpropagation computes gradients; an \textbf{optimiser} uses them to update weights\textsuperscript{\cite[][pp.~22--25]{fleuretLittleBookDeep}}.
    \tcblower
    \begin{tblr}{|Q[l,m]|X[l]|}
        \hline
        \textbf{Optimiser} & \textbf{Update Rule / Key Idea} \\
        \hline
        SGD & $\theta \leftarrow \theta - \eta \nabla\mathcal{L}$ \\
        \hline
        SGD + Momentum & Accumulates past gradients; smooths updates \\
        \hline
        RMSprop & Per-parameter adaptive learning rate \\
        \hline
        Adam & Momentum + RMSprop; default choice\textsuperscript{\cite[][]{kingmaAdamMethodStochastic2017}} \\
        \hline
    \end{tblr}
\captionof{table}{Backpropagation and Gradient Descent: Gradient Descent Variants.}
\label{tab:ch06-deep-learning-inline-04}

    \tcbitem[title={Mini-Batch Gradient Descent}, raster multicolumn=6]
    In practice, gradients are computed on \textbf{mini-batches} (subsets of the training data, typically 32--512 examples)\textsuperscript{\cite[][p.~23]{fleuretLittleBookDeep}}.
    This provides a compromise between full-batch gradient descent (low noise, expensive) and single-sample SGD (high noise, cheap).
    Mini-batch SGD introduces beneficial stochastic noise that helps escape shallow local minima and saddle points.
    The \textbf{learning rate} $\eta$ is the most important hyperparameter: too large $\to$ divergence; too small $\to$ slow convergence.
\end{tcbitemize}

% ---- 6.3.4 Training Loop -------------------------------------------------------------------
\subsection{Training Loop}

\begin{tcbitemize}[ skin=sectionraster ]
    \tcbitem[title={Epochs and Batches}, raster multicolumn=3]
    The training process iterates through the dataset in structured passes\textsuperscript{\cite[][pp.~24--26]{fleuretLittleBookDeep}}:
    \begin{itemize}
        \item \textbf{Epoch}: one complete pass through the entire training set
        \item \textbf{Iteration/step}: one parameter update using a single mini-batch
        \item \textbf{Batch size}: number of examples per mini-batch; affects gradient quality and memory usage
    \end{itemize}
    A typical training run uses tens to hundreds of epochs, with the dataset reshuffled between epochs.

    \tcbitem[title={Data Splits and Early Stopping}, raster multicolumn=3]
    The data is split into \textbf{training}, \textbf{validation}, and \textbf{test} sets (common split: 70/15/15 or 80/10/10)\textsuperscript{\cite[][p.~26]{fleuretLittleBookDeep}}.
    During training, the model is evaluated on the validation set after each epoch.
    \textbf{Early stopping}: if validation loss does not improve for a set number of epochs (``patience''), training halts and the model reverts to the best checkpoint.
    This is one of the most effective and simplest forms of regularisation.
    \tcblower
    \textbf{Learning rate scheduling} further improves training:
    \begin{itemize}
        \item \emph{Step decay}: reduce $\eta$ by a factor at fixed intervals
        \item \emph{Cosine annealing}: smoothly decrease $\eta$ following a cosine curve
        \item \emph{Warmup}: start with a very small $\eta$ and linearly increase it over the first few epochs
    \end{itemize}
\end{tcbitemize}

% ---- 6.3.5 Regularization and Overtraining --------------------------------------------------
\subsection{Regularization and Overtraining}

\begin{tcbitemize}[ skin=sectionraster ]
    \tcbitem[title={Overfitting}, raster multicolumn=6]
    A model \textbf{overfits} when it memorises the training data (low training loss) but generalises poorly to unseen data (high validation loss)\textsuperscript{\cite[][pp.~27--29]{fleuretLittleBookDeep}}.
    Deep networks with millions of parameters are particularly prone to overfitting when training data is limited.
    \textbf{Regularisation} techniques constrain the model's capacity to reduce overfitting.

    \tcbitem[title={Dropout}, raster multicolumn=3]
    During each training step, \textbf{dropout} randomly sets a fraction $p$ of neuron activations to zero\textsuperscript{\cite[][p.~28]{fleuretLittleBookDeep}}.
    This prevents co-adaptation: neurons cannot rely on specific other neurons being active, forcing redundant representations.
    At test time, dropout is turned off and activations are scaled by $(1 - p)$ to maintain expected values.

    \tcbitem[title={Weight Regularisation}, raster multicolumn=3]
    Adding a penalty term to the loss discourages large weights:
    \begin{itemize}
        \item \textbf{L2} (weight decay): $\mathcal{L}_{\mathrm{reg}} = \mathcal{L} + \lambda \sum w_i^2$ --- encourages small, distributed weights
        \item \textbf{L1}: $\mathcal{L}_{\mathrm{reg}} = \mathcal{L} + \lambda \sum |w_i|$ --- encourages sparsity (many weights exactly zero)
    \end{itemize}
    In practice, L2 regularisation (weight decay) is the most commonly used, often built directly into the optimiser.

    \tcbitem[title={Batch Normalization}, raster multicolumn=3]
    \textbf{Batch Normalization} (BN)\textsuperscript{\cite[][]{ioffeBatchNormalizationAccelerating2015}} normalises the activations of each layer across the mini-batch to have zero mean and unit variance, then applies learned scale $\gamma$ and shift $\beta$ parameters.
    Benefits: stabilises training, allows higher learning rates, and provides a mild regularisation effect.
    BN is applied between the linear transformation and the activation function.

    \tcbitem[title={Data Augmentation}, raster multicolumn=3]
    \textbf{Data augmentation} artificially expands the training set by applying label-preserving transformations\textsuperscript{\cite[][p.~29]{fleuretLittleBookDeep}}:
    random flips, rotations, crops, colour jitter, or more advanced techniques like mixup and cutout.
    This is particularly effective for image tasks and can reduce overfitting significantly without changing the model architecture.
\end{tcbitemize}


% ==== Section 4: Alternative Training Methods ================================================
\section{Alternative Training Methods}

% ---- 6.4.1 Attention -----------------------------------------------------------------------
\subsection{Attention}

\begin{tcbitemize}[ skin=sectionraster ]
    \tcbitem[title={Attention: Why It Matters}, raster multicolumn=3]
    Fully connected layers cannot handle variable-length inputs; convolutional layers propagate information only locally.
    \textbf{Attention layers} address this by computing, for every component of the output, a weighted average over the \emph{entire} input tensor --- regardless of distance\textsuperscript{\cite[][p.~73]{fleuretLittleBookDeep}}.
    This makes them the key building block for Transformers and modern large language models.

    \tcbitem[title={Query, Key, Value}, raster multicolumn=3]
    The \textbf{attention operator} att$(Q, K, V)$ takes three tensors\textsuperscript{\cite[][pp.~74--76]{fleuretLittleBookDeep}}:
    \begin{itemize}
        \item $Q \in \mathbb{R}^{N^q \times D^{qk}}$ --- \emph{Queries}
        \item $K \in \mathbb{R}^{N^{kv} \times D^{qk}}$ --- \emph{Keys}
        \item $V \in \mathbb{R}^{N^{kv} \times D^v}$ --- \emph{Values}
    \end{itemize}
    For each query $Q_n$, attention scores $A_{n,m}$ measure how well $Q_n$ matches each key $K_m$.
    The retrieved value $Y_n = \sum_m A_{n,m} V_m$ is a weighted average of all values.

    \tcbitem[title={Scaled Dot-Product Attention}, raster multicolumn=6]
    Attention scores are computed as a scaled softmax of dot-products between queries and keys\textsuperscript{\cite[][p.~3]{vaswaniAttentionAllYou2023}}.
    The scaling factor $1/\!\sqrt{D^{qk}}$ keeps the dot products in a range where softmax has useful gradients.
    \tcblower
    \begin{equation}
        \mathrm{Attention}(Q, K, V)
        = \mathrm{softmax}\!\left(\frac{QK^{\!\top}}{\sqrt{D^{qk}}}\right) V
    \end{equation}

    \tcbitem[title={Self- and Cross-Attention}, raster multicolumn=3]
    In a \textbf{self-attention} block, the three input sequences $X^q$, $X^k$, $X^v$ are identical\textsuperscript{\cite[][p.~79]{fleuretLittleBookDeep}}.
    Every position can gather information from every other position in the same sequence --- enabling global context.

    In a \textbf{cross-attention} block, $X^k$ and $X^v$ come from a different sequence (e.g., the encoder representation), while $X^q$ comes from the current sequence (e.g., the decoder).

    \tcbitem[title={Multi-Head Attention}, raster multicolumn=3]
    A single attention head averages over all value positions and cannot attend to multiple subspaces simultaneously.
    \textbf{Multi-Head Attention}\textsuperscript{\cite[][pp.~4--5]{vaswaniAttentionAllYou2023}} runs $H$ attention heads in parallel, each with its own learned projections $W^q_h, W^k_h, W^v_h$, then concatenates and projects:
    \begin{equation}
        \mathrm{MultiHead}(Q,K,V) = \mathrm{Concat}(\mathrm{head}_1,\ldots,\mathrm{head}_H)\,W^O
    \end{equation}
    Each head can learn to attend to a different type of relationship (syntactic, semantic, positional).
\end{tcbitemize}

% ---- 6.4.2 Transformer Architecture ---------------------------------------------------------
\subsection{Transformer Architecture}

\begin{tcbitemize}[ skin=sectionraster ]
    \tcbitem[title={Why Transformers Replaced RNNs}, raster multicolumn=3]
    Recurrent networks process sequences step-by-step, requiring $O(n)$ sequential operations --- preventing parallelism and making long-range dependencies hard to learn\textsuperscript{\cite[][p.~6]{vaswaniAttentionAllYou2023}}.
    Self-attention layers connect all positions with $O(1)$ sequential operations and constant path length between any two positions.
    Transformers train significantly faster on modern hardware and scale to billions of parameters\textsuperscript{\cite[][p.~93]{fleuretLittleBookDeep}}.

    \tcbitem[title={Encoder--Decoder Structure}, raster multicolumn=3]
    The original Transformer\textsuperscript{\cite[][p.~2]{vaswaniAttentionAllYou2023}} is designed for sequence-to-sequence tasks (e.g., translation):
    \begin{itemize}
        \item \textbf{Encoder:} $N$ identical layers, each with (1) multi-head self-attention and (2) position-wise feed-forward network (FFN); residual connections and layer normalisation around each sub-layer
        \item \textbf{Decoder:} $N$ layers with (1) masked causal self-attention, (2) cross-attention over encoder output, and (3) FFN; causal masking prevents attending to future tokens
    \end{itemize}

    \tcbitem[title={Transformer Block Diagram}, raster multicolumn=6]
    The encoder processes the full input in parallel; the decoder generates output tokens autoregressively.
    \tcblower
    \begin{center}
    \begin{tikzpicture}[
        node distance=0.25cm and 0.7cm,
        box/.style={draw, rounded corners=2pt, minimum width=2.8cm, minimum height=0.55cm,
                    font=\sffamily\scriptsize, align=center, fill=blue!10},
        addnorm/.style={draw, rounded corners=2pt, minimum width=2.8cm, minimum height=0.45cm,
                    font=\sffamily\scriptsize, align=center, fill=orange!15},
        lbl/.style={font=\sffamily\scriptsize, text=gray},
        arr/.style={-{Stealth[length=2mm]}, thick},
        brace/.style={decorate, decoration={brace, amplitude=4pt}}
    ]
        % ---- ENCODER (left column) ----
        \node[box] (emb_e)                          {Embedding};
        \node[box, above=of emb_e] (pe_e)            {+ Pos.\ Encoding};
        \node[addnorm, above=of pe_e] (an1_e)       {Add \& Norm};
        \node[box, above=of an1_e] (mha_e)          {Multi-Head\\Self-Attention};
        \node[addnorm, above=of mha_e] (an2_e)      {Add \& Norm};
        \node[box, above=of an2_e] (ffn_e)          {Feed Forward};
        \node[addnorm, above=of ffn_e] (an3_e)      {Add \& Norm};
        \node[lbl, left=0.15cm of mha_e] {$\times N$};

        % ---- DECODER (right column) ----
        \node[box, right=2.5cm of emb_e] (emb_d)         {Embedding};
        \node[box, above=of emb_d] (pe_d)           {+ Pos.\ Encoding};
        \node[addnorm, above=of pe_d] (an1_d)      {Add \& Norm};
        \node[box, above=of an1_d] (cmha_d)        {Masked\\Self-Attention};
        \node[addnorm, above=of cmha_d] (an2_d)    {Add \& Norm};
        \node[box, above=of an2_d] (xatt_d)        {Cross-Attention};
        \node[addnorm, above=of xatt_d] (an3_d)    {Add \& Norm};
        \node[box, above=of an3_d] (ffn_d)         {Feed Forward};
        \node[addnorm, above=of ffn_d] (an4_d)     {Add \& Norm};
        \node[lbl, right=0.15cm of cmha_d] {$\times N$};

        % Output
        \node[box, above=0.3cm of an4_d, fill=green!15] (out) {Linear + Softmax};

        % Encoder arrows
        \draw[arr] (emb_e) -- (pe_e);
        \draw[arr] (pe_e) -- (an1_e);
        \draw[arr] (an1_e) -- (mha_e);
        \draw[arr] (mha_e) -- (an2_e);
        \draw[arr] (an2_e) -- (ffn_e);
        \draw[arr] (ffn_e) -- (an3_e);

        % Decoder arrows
        \draw[arr] (emb_d) -- (pe_d);
        \draw[arr] (pe_d) -- (an1_d);
        \draw[arr] (an1_d) -- (cmha_d);
        \draw[arr] (cmha_d) -- (an2_d);
        \draw[arr] (an2_d) -- (xatt_d);
        \draw[arr] (xatt_d) -- (an3_d);
        \draw[arr] (an3_d) -- (ffn_d);
        \draw[arr] (ffn_d) -- (an4_d);
        \draw[arr] (an4_d) -- (out);

        % Cross-attention connection from encoder
        \draw[arr, dashed, gray] (an3_e.east) .. controls +(0.4,0.8) and +(-0.4,0) .. (xatt_d.west)
            node[midway, above, font=\sffamily\tiny, gray] {K, V};

        % Input labels
        \node[lbl, below=0.1cm of emb_e] {Source tokens};
        \node[lbl, below=0.1cm of emb_d] {Target tokens (shifted)};
        \node[lbl, above=0.1cm of out] {Output probs};
    \end{tikzpicture}
\captionof{figure}{Transformer Architecture: Transformer Block Diagram.}
\label{fig:ch06-deep-learning-inline-02}
    \end{center}

    \tcbitem[title={Positional Encoding}, raster multicolumn=3]
    Attention is permutation-invariant --- it has no notion of sequence order.
    Position information is injected by adding a \textbf{positional encoding} to the input embeddings\textsuperscript{\cite[][pp.~81--82]{fleuretLittleBookDeep}}.
    The original Transformer uses sinusoidal encoding:
    \begin{equation}
        \mathrm{PE}_{t,2i} = \sin\!\left(\frac{t}{10^{4 \cdot 2i/D}}\right),
        \quad
        \mathrm{PE}_{t,2i+1} = \cos\!\left(\frac{t}{10^{4 \cdot 2i/D}}\right)
    \end{equation}
    The deterministic encoding can be evaluated at positions beyond those observed during training, but useful length extrapolation is an empirical property rather than a mathematical guarantee. Learned positional embeddings performed similarly in the original experiments\textsuperscript{\cite[][p.~6]{vaswaniAttentionAllYou2023}}.

    \tcbitem[title={GPT and BERT: Decoder- and Encoder-Only}, raster multicolumn=3]
    Modern large language models are typically \emph{decoder-only} or \emph{encoder-only}\textsuperscript{\cite[][pp.~97--99]{fleuretLittleBookDeep}}:
    \begin{itemize}
        \item \textbf{GPT} (Radford et al.\ 2018): decoder-only; stacked causal self-attention blocks; trained autoregressively on next-token prediction; scales to hundreds of billions of parameters
        \item \textbf{BERT}\textsuperscript{\cite[][]{devlinBERTPretrainingDeep2019}} (Devlin et al.\ 2019): encoder-only; bidirectional; trained on masked language modeling and next-sentence prediction
    \end{itemize}
\end{tcbitemize}

% ---- 6.4.3 Feedback Alignment ---------------------------------------------------------------
\subsection{Feedback Alignment}

\begin{tcbitemize}[ skin=sectionraster ]
    \tcbitem[title={The Weight Transport Problem}, raster multicolumn=3]
    Standard backpropagation transmits error signals using the \emph{transpose} of the forward weight matrix $W^\top$\textsuperscript{\cite[][p.~3]{fischerVorlesungGrundlagenNeuronale2024}}.
    This requires the backward path to have exact knowledge of every forward weight --- a condition with no known biological mechanism, called the \textbf{weight transport problem}.

    \tcbitem[title={Feedback Alignment}, raster multicolumn=3]
    Lillicrap et al.\ (2016) showed that replacing $W^\top$ with a \textbf{fixed random feedback matrix} $B$ still produces learning.
    The forward weights $W$ gradually align with $B$ during training --- the network effectively ``shapes itself'' to match the random feedback.
    \textbf{Direct Feedback Alignment} (N{\o}kland, 2016) extends this further: each layer receives error signals \emph{directly} from the output layer via separate random matrices, bypassing the chain entirely.
\end{tcbitemize}

% ---- 6.4.4 Synthetic Gradients --------------------------------------------------------------
\subsection{Synthetic Gradients}

\begin{tcbitemize}[ skin=sectionraster ]
    \tcbitem[title=Decoupled Neural Interfaces, raster multicolumn=6]
    In standard backpropagation, no layer can update until the full forward and backward passes complete --- a constraint called \textbf{update locking}.
    Jaderberg et al.\ (2017) introduced \textbf{Decoupled Neural Interfaces (DNI)}: a small auxiliary network $M_i$ attached to each layer that \emph{predicts} the gradient $\partial\mathcal{L}/\partial h_i$ without waiting for it to arrive.
    The layer updates immediately using this synthetic gradient; $M_i$ is itself trained by the true gradient when it eventually arrives.
    This enables asynchronous, modular training of deep networks --- each layer acts as an independent learning module with a local objective.
    \tcblower
    \begin{center}
    \begin{tikzpicture}[
        node distance=0.35cm and 0.5cm,
        lay/.style={draw, rounded corners, minimum width=1.8cm, minimum height=0.6cm,
                    font=\sffamily\scriptsize, fill=blue!10},
        dni/.style={draw, rounded corners, minimum width=1.4cm, minimum height=0.5cm,
                    font=\sffamily\scriptsize, fill=orange!20},
        arr/.style={-{Stealth[length=2mm]}, thick},
        darr/.style={-{Stealth[length=2mm]}, dashed, gray}
    ]
        \node[lay] (L1) {Layer 1};
        \node[lay, right=of L1] (L2) {Layer 2};
        \node[lay, right=of L2] (L3) {Layer 3};
        \node[lay, right=of L3] (loss) {Loss};

        \node[dni, below=of L1] (M1) {DNI $M_1$};
        \node[dni, below=of L2] (M2) {DNI $M_2$};

        \draw[arr] (L1) -- (L2) node[midway, above, font=\sffamily\tiny] {$h_1$};
        \draw[arr] (L2) -- (L3) node[midway, above, font=\sffamily\tiny] {$h_2$};
        \draw[arr] (L3) -- (loss);

        \draw[arr] (L1) -- (M1);
        \draw[arr] (L2) -- (M2);

        \draw[darr] (M1) -- (L1) node[midway, right, font=\sffamily\tiny, gray] {$\tilde{\delta}_1$};
        \draw[darr] (M2) -- (L2) node[midway, right, font=\sffamily\tiny, gray] {$\tilde{\delta}_2$};
    \end{tikzpicture}
\captionof{figure}{Synthetic Gradients: Decoupled Neural Interfaces.}
\label{fig:ch06-deep-learning-inline-03}
    \end{center}
\end{tcbitemize}

% ---- 6.4.5 Decoupled Network Interfaces -----------------------------------------------------
\subsection{Decoupled Network Interfaces}

\begin{tcbitemize}[ skin=sectionraster ]
    \tcbitem[title={Local Learning Signals}, raster multicolumn=6]
    The full Decoupled Neural Interface framework removes both \emph{update locking} (no backward pass needed) and \emph{forward locking} (synthetic activations can substitute for real forward signals).
    Each module then has a \textbf{local learning signal} and can train asynchronously.
    This is more aligned with how biological neural circuits are thought to operate --- local Hebbian-like updates rather than a global error signal broadcast through the entire network\textsuperscript{\cite[][p.~2]{fischerVorlesungGrundlagenNeuronale2024}}.
    In practice, synthetic gradients are primarily a research contribution; standard backpropagation remains the dominant approach.
\end{tcbitemize}


% ==== Section 5: Further Network Architectures ===============================================
\section{Further Network Architectures}

% ---- 6.5.1 GANs -----------------------------------------------------------------------------
\subsection{Generative Adversarial Networks}

\begin{tcbitemize}[ skin=sectionraster ]
    \tcbitem[title={Generator and Discriminator}, raster multicolumn=3]
    A \textbf{Generative Adversarial Network (GAN)} introduced by Goodfellow et al.\ (2014) pits two networks against each other\textsuperscript{\cite[][p.~125]{fleuretLittleBookDeep}}:
    \begin{itemize}
        \item \textbf{Generator} $G(z)$: takes random noise $z \sim p(z)$ and produces a structured output (e.g., an image) that mimics real data
        \item \textbf{Discriminator} $D(x)$: takes either a real sample or $G(z)$; outputs probability that the input is real
    \end{itemize}
    At equilibrium, $G$ produces samples indistinguishable from real data and $D$ outputs $0.5$ everywhere.

    \tcbitem[title={The Minimax Objective}, raster multicolumn=3]
    Training is formulated as a two-player minimax game\textsuperscript{\cite[][p.~125]{fleuretLittleBookDeep}}: $D$ maximises its ability to distinguish real from fake; $G$ minimises the probability that $D$ identifies its outputs as fake.
    \tcblower
    \begin{equation}
        \min_G \max_D \;\bigl[\,\mathbb{E}[\log D(x)] + \mathbb{E}[\log(1 - D(G(z)))]\,\bigr]
    \end{equation}

    \tcbitem[title={Training Dynamics and Mode Collapse}, raster multicolumn=6]
    GAN training is notoriously unstable.
    When the discriminator is too strong, gradients to the generator vanish; when it is too weak, the generator does not learn useful structure.
    A key failure mode is \textbf{mode collapse}: the generator learns to produce a small variety of outputs that fool the discriminator, ignoring most of the real data distribution.
    As gradients flow from $D$ back to $G$, the discriminator informs the generator about the cues it detects --- the generator then learns to eliminate those cues\textsuperscript{\cite[][p.~125]{fleuretLittleBookDeep}}.
    Variants such as Wasserstein GAN (WGAN) and spectral normalisation address training stability.
\end{tcbitemize}

% ---- 6.5.2 Autoencoders ---------------------------------------------------------------------
\subsection{Autoencoders}

\begin{tcbitemize}[ skin=sectionraster ]
    \tcbitem[title={Encoder--Bottleneck--Decoder}, raster multicolumn=3]
    An \textbf{autoencoder} consists of an \emph{encoder} $f_\phi$ and a \emph{decoder} $g_\theta$\textsuperscript{\cite[][p.~125]{fleuretLittleBookDeep}}.
    The encoder maps the high-dimensional input $x$ to a low-dimensional latent vector $z = f_\phi(x)$; the decoder reconstructs $\hat{x} = g_\theta(z)$.
    The bottleneck forces the network to learn a compact representation of the data manifold.
    \tcblower
    \begin{itemize}
        \item Loss: $\mathcal{L} = \|x - \hat{x}\|^2$ (MSE) or cross-entropy
        \item Applications: dimensionality reduction, denoising, anomaly detection, pre-training
        \item A \textbf{denoising autoencoder} receives corrupted $\tilde{x}$ as input but targets clean $x$, forcing more robust representations
    \end{itemize}

    \tcbitem[title={Variational Autoencoder (VAE)}, raster multicolumn=3]
    The \textbf{Variational Autoencoder} (Kingma \& Welling, 2013) imposes a prior distribution on the latent space\textsuperscript{\cite[][p.~125]{fleuretLittleBookDeep}}.
    The encoder outputs parameters $(\mu, \sigma)$ of a Gaussian; the \emph{reparameterisation trick} ($z = \mu + \sigma \odot \varepsilon$, $\varepsilon \sim \mathcal{N}(0,I)$) makes sampling differentiable.
    \tcblower
    \begin{equation}
        \mathcal{L}_{\mathrm{ELBO}} = \mathbb{E}_{q_\phi(z|x)}\bigl[\log p_\theta(x|z)\bigr]
        - \mathrm{KL}\bigl(q_\phi(z|x) \,\|\, p(z)\bigr)
    \end{equation}
    The KL term regularises the latent space towards $\mathcal{N}(0,I)$, enabling new samples to be drawn by sampling $z \sim \mathcal{N}(0,I)$ and decoding.
\end{tcbitemize}

% ---- 6.5.3 Restricted Boltzmann Machines ----------------------------------------------------
\subsection{Restricted Boltzmann Machines}

\begin{tcbitemize}[ skin=sectionraster ]
    \tcbitem[title={Energy-Based Generative Model}, raster multicolumn=3]
    A \textbf{Restricted Boltzmann Machine (RBM)} is a bipartite undirected graphical model with visible units $v$ and hidden units $h$, and no intra-layer connections.
    The joint probability is defined via an energy function:
    \begin{equation}
        E(v, h) = -v^\top W h - b^\top v - c^\top h
    \end{equation}
    with $P(v, h) \propto \exp(-E(v,h))$.
    The restriction to bipartite structure makes conditional distributions tractable: $P(h|v)$ and $P(v|h)$ are independent Bernoulli.

    \tcbitem[title={Contrastive Divergence Training}, raster multicolumn=3]
    Exact maximum-likelihood training requires computing the partition function (intractable).
    \textbf{Contrastive Divergence (CD-k)} approximates the gradient\textsuperscript{\cite[][]{Kriesel2007NeuralNetworks}}:
    \begin{itemize}
        \item \emph{Positive phase}: clamp real data $v$; sample $h \sim P(h|v)$
        \item \emph{Negative phase}: run $k$ steps of Gibbs sampling to get model sample $(\tilde{v},\tilde{h})$
        \item Update: $\Delta W \propto \langle vh \rangle_\mathrm{data} - \langle \tilde{v}\tilde{h} \rangle_\mathrm{model}$
    \end{itemize}
    RBMs were used for greedy layer-wise pre-training of deep belief networks (Hinton et al., 2006) but have been largely superseded by modern training methods.
\end{tcbitemize}

% ---- 6.5.4 Capsule Networks -----------------------------------------------------------------
\subsection{Capsule Networks}

\begin{tcbitemize}[ skin=sectionraster ]
    \tcbitem[title={Capsules and Dynamic Routing}, raster multicolumn=6]
    Proposed by Sabour, Frosst \& Hinton (2017) to address a fundamental limitation of CNNs: pooling operations are viewpoint-invariant but discard spatial and pose information.
    A \textbf{capsule} is a group of neurons whose \emph{activity vector} encodes both the presence (vector length) and instantiation parameters such as pose, scale, and orientation (vector direction) of an entity.
    Lower-level capsules send \textbf{votes} to candidate higher-level capsules; an iterative \textbf{dynamic routing} algorithm (3 iterations) selects the higher-level capsule whose pose best agrees with the votes.
    This enables \textbf{equivariance}: as viewpoint changes, the capsule representation changes predictably, preserving part-whole relationships\textsuperscript{\cite[][]{fleuretLittleBookDeep}}.
    \tcblower
    \begin{itemize}
        \item Advantage: encodes spatial hierarchies, more interpretable than CNNs
        \item Limitation: $O(n^2)$ routing cost; does not scale well to complex datasets such as ImageNet
    \end{itemize}
\end{tcbitemize}

% ---- 6.5.5 Spiking Networks -----------------------------------------------------------------
\subsection{Spiking Networks}

\begin{tcbitemize}[ skin=sectionraster ]
    \tcbitem[title={Third-Generation Neural Networks}, raster multicolumn=6]
    \textbf{Spiking Neural Networks (SNNs)}, introduced by Maass (1997), are described as the ``third generation'' of neural networks.
    Rather than communicating through continuous activations, neurons emit discrete \textbf{spike} events.
    The \textbf{Leaky Integrate-and-Fire (LIF)} neuron model accumulates input in a membrane potential $V(t)$; when $V$ exceeds a threshold $\theta$, a spike is emitted and $V$ resets.
    Information can be encoded in spike rates or precise spike timing.
    \tcblower
    \begin{itemize}
        \item \textbf{Energy efficiency}: sparse spike events map naturally onto neuromorphic hardware (Intel Loihi, IBM TrueNorth) with orders-of-magnitude lower power than GPU-based ANNs
        \item \textbf{Biological plausibility}: closer to cortical computation than rate-coded networks
        \item \textbf{Training challenge}: the spike function is non-differentiable; \emph{surrogate gradients} (smooth approximations) enable gradient-based training
        \item \textbf{Current status}: promising for edge and embedded AI; not yet competitive with ANNs on standard benchmarks
    \end{itemize}
\end{tcbitemize}


% ==== Section 6: Further Reading =============================================================
\section{Further Reading}

\begin{tcbitemize}[ skin=sectionraster ]
    \tcbitem[title=Recommended Sources for Deep Learning, raster multicolumn=6]
    \begin{itemize}
        \item \fullcite{fleuretLittleBookDeep} — Concise yet comprehensive overview of deep learning: architectures, training, attention, transformers, and generative models.
        \item \fullcite{Kriesel2007NeuralNetworks} — Thorough introduction to biological foundations, perceptrons, MLPs, and RBMs with historical context.
        \item \fullcite{robertsPrinciplesDeepLearning2022} — Theoretical perspective on initialization, depth, and why deep networks generalize.
        \item \fullcite{vaswaniAttentionAllYou2023} — The original Transformer paper: self-attention, multi-head attention, positional encoding, and encoder--decoder architecture.
        \item \fullcite{devlinBERTPretrainingDeep2019} — BERT: bidirectional pre-training with masked language modelling and next-sentence prediction.
        \item \fullcite{fischerVorlesungGrundlagenNeuronale2024} — German-language lecture notes covering neural network foundations, feedback alignment, and biological plausibility.
        \item \fullcite{dhamaniIntroductionGenerativeAI2024} — Practical introduction to generative AI: LLMs, GANs, VAEs, and transformer-based generation.
    \end{itemize}
\end{tcbitemize}
